\section{Choice of Output: MIDI as a Semantic Target}\label{sec:midi-target}

The final component of the methodology is the selection of the output format.
While some musical DSLs target raw audio synthesis or real-time event streaming via Open Sound Control (OSC),
this research utilizes the \textit{MIDI Format 1} protocol as its primary build target.

\subsection{Interoperability and the Producer Workflow}

A key goal of this DSL is to serve as a tool for \textit{producers} and \textit{composers}, not just a
self-contained instrument.
By targeting MIDI, the compiler produces an artifact that is natively understood by almost every Digital
Audio Workstation (DAW) and notation software in existence.
This interoperability allows the "code-based composition" to be seamlessly integrated into a traditional
production workflow, where the resulting MIDI can drive high-end virtual instruments, orchestral libraries,
or external hardware synthesizers.

\subsection{Preserving Symbolic Identity}

MIDI is a \textit{symbolic representation} of music.
Unlike audio, which is a stream of samples, MIDI preserves the identity of notes, velocities, and rhythmic placement.
This makes MIDI an ideal target for a structured DSL because the high-level constructs of the language (e.g.,
a "C4" note at beat 2) map directly to low-level MIDI messages (\texttt{Note-On, key 72, velocity 100}).
This one-to-one mapping ensures that the semantic intent of the composer is perfectly preserved in the output
file, allowing for further offline analysis or modification in a DAW.

\subsection{Offline Analysis and Reproducibility}

Finally, targeting a static file format like MIDI ensures absolute \textit{reproducibility}.
In a live coding environment, the same script might sound different on different machines or under different CPU loads.
With the compiled approach, the compiler generates a deterministic binary file.
This file can be archived, version-controlled, and shared, guaranteeing that the composition remains "frozen"
exactly as the composer intended, regardless of the playback environment.
This methodology reinforces the thesis that musical composition is a form of structured data that benefits
from the same rigorous processing as software source code.
